
% =====================================================================
% 1. Introduction
% =====================================================================
\section{Introduction}
\label{sec:intro}

A weather station can produce thousands of readings per day. A network of
hundreds of stations produces more data than a person can inspect manually.
Yet those measurements feed flood warnings, climate records, research and
operational decisions. Somewhere between the sensor and the user, software
must decide which readings look trustworthy and which deserve attention.

That decision is harder than it first appears. A value may be physically
impossible, may jump too quickly, may stop changing, may disagree with a
neighbour, or may simply be an unusual but real event. A useful QC system
must detect faults without turning every rare event into an error. It must
also explain why it objected, because a flag is ultimately read by a person.

My interest in this problem came from an internship in which I extended and
audited an automated QC pipeline. The code contained familiar checks and
many configuration constants. That raised questions the code itself could
not answer: Why this statistic? Why this threshold? Which assumptions make
it valid? What happens when the same method is moved from an independent
sample to a correlated sensor stream?

This report answers those questions by reading five works in chronological
order and then testing the main quantitative ideas on a benchmark with known
faults. The five core readings are listed in Table~\ref{tab:fiveworks}.

\begin{table}[h]
\centering
\small
\caption{The five core readings and the role each plays in the report.}
\label{tab:fiveworks}
\begin{tabular}{@{}p{3.5cm}p{1.2cm}p{8.7cm}@{}}
\toprule
Work & Year & Why it matters \\
\midrule
\citet{grubbs1969} & 1969 & Classical hypothesis-test view of an outlying observation. \\
\citet{iglewicz1993} & 1993 & Robust modified z-score using the median and MAD; source of the 3.5 screening rule. \\
\citet{shafer2000} & 2000 & A complete operational QC system for a large automatic weather network. \\
\citet{campbell2013} & 2013 & General framework for QA/QC in streaming environmental sensor data. \\
\citet{leigh2019} & 2019 & Comparison of rule-based, regression and feature-based anomaly detection on labelled water-quality data. \\
\bottomrule
\end{tabular}
\end{table}

The course suggests five roles: Historian, Technician, Experimentalist,
Futurist and Critic. I use one section for each role. This keeps the report
close to the course structure and makes the main argument easy to follow.

% =====================================================================
% 2. Historian
% =====================================================================
